%! Inverse Trigonometric Functions — Unit 7, Lesson 7.2 — Guided Notes (Answer Key)
\documentclass[10pt]{article}
\usepackage{precalculus-article}
\usepackage{precalculus-key}   % pulls in -boxes; do NOT also load -boxes
\newcommand{\TallMath}[1]{$\displaystyle #1\rule[-1.4em]{0pt}{3.2em}$}

\begin{document}

\pageheader{Unit 7, Lesson 7.2}{Guided Notes --- Answer Key}

\vspace{0.10in}

%----------------------------------------------------------------------
\begin{notesbox}{LO 3.9.A \; Inverse Trigonometric Functions --- Key}

%--- Vocabulary row
\begin{skillbox}[Key Vocabulary]{greenbox}
\begin{multicols}{2}
\begin{itemize}[leftmargin=1.4em, itemsep=1pt]
    \item \textbf{arcsine} --- the inverse of the restricted sine function;
          outputs an angle in $[-\tfrac{\pi}{2},\tfrac{\pi}{2}]$
    \item \textbf{arccosine} --- the inverse of the restricted cosine function;
          outputs an angle in $[0,\pi]$
    \item \textbf{arctangent} --- the inverse of the restricted tangent function;
          outputs an angle in $(-\tfrac{\pi}{2},\tfrac{\pi}{2})$
    \item \textbf{restricted domain} --- a subset of the natural domain chosen
          so a function is one-to-one
    \item \textbf{principal value} --- the unique output angle in the restricted
          range of an inverse trig function
\end{itemize}
\end{multicols}
\end{skillbox}

\vspace{0.08in}

\textbf{Section 1 --- Why Do We Need Restricted Domains?}

The table below shows selected values of $f(x) = \sin x$.

\begin{center}
\renewcommand{\arraystretch}{1.8}
\begin{tabular}{|c|c|c|c|c|c|c|c|}
\hline
$x$ &
$-\dfrac{\pi}{2}$ & $-\dfrac{\pi}{6}$ & $0$ &
$\dfrac{\pi}{6}$ & $\dfrac{\pi}{2}$ & $\dfrac{5\pi}{6}$ & $\pi$ \\
\hline
$\sin x$ & $-1$ & $-\tfrac{1}{2}$ & $0$ &
$\tfrac{1}{2}$ & $1$ & $\tfrac{1}{2}$ & $0$ \\
\hline
\end{tabular}
\end{center}

$\sin x$ is \ans{NOT one-to-one}.

Because the sine function is periodic, it \ans{fails} the horizontal line test,
so it \ans{does not} have an inverse over its full domain.

\vspace{0.08in}

\textbf{Section 2 --- Restricted Domains and the Inverse Trig Functions}

\begin{center}
\renewcommand{\arraystretch}{2.2}
\begin{tabular}{|l|c|c|c|c|}
\hline
\textbf{Function} & \textbf{Restricted Domain} & \textbf{Range} &
\textbf{Inverse Name} & \textbf{Notation} \\
\hline
Sine    & \ans{$\left[-\dfrac{\pi}{2},\dfrac{\pi}{2}\right]$} & $[-1,\,1]$  & Arcsine   & $\arcsin x$ \\
\hline
Cosine  & \ans{$[0,\,\pi]$}   & $[-1,\,1]$         & Arccosine & $\arccos x$ \\
\hline
Tangent & \ans{$\left(-\dfrac{\pi}{2},\dfrac{\pi}{2}\right)$} & $(-\infty,\infty)$ & Arctangent & $\arctan x$ \\
\hline
\end{tabular}
\end{center}

The output of an inverse trig function is an \ans{angle} measure (the principal value).

\vspace{0.08in}

\textbf{Section 3 --- Evaluating Inverse Trig Functions}

\begin{center}
\renewcommand{\arraystretch}{2.4}
\begin{tabular}{|l|l|l|}
\hline
\textbf{Expression} & \textbf{Think:} & \textbf{Answer} \\
\hline
$\arcsin\!\left(\dfrac{1}{2}\right)$ &
  Which $\theta\in[-\tfrac{\pi}{2},\tfrac{\pi}{2}]$ has $\sin\theta=\tfrac{1}{2}$? &
  \ans{$\dfrac{\pi}{6}$} \\
\hline
$\arccos(0)$ &
  Which $\theta\in[0,\pi]$ has $\cos\theta=0$? &
  \ans{$\dfrac{\pi}{2}$} \\
\hline
$\arctan(1)$ &
  Which $\theta\in(-\tfrac{\pi}{2},\tfrac{\pi}{2})$ has $\tan\theta=1$? &
  \ans{$\dfrac{\pi}{4}$} \\
\hline
$\arcsin\!\left(-\dfrac{\sqrt{2}}{2}\right)$ &
  Which $\theta\in[-\tfrac{\pi}{2},\tfrac{\pi}{2}]$ has $\sin\theta=-\dfrac{\sqrt{2}}{2}$? &
  \ans{$-\dfrac{\pi}{4}$} \\
\hline
$\arccos(-1)$ &
  Which $\theta\in[0,\pi]$ has $\cos\theta=-1$? &
  \ans{$\pi$} \\
\hline
\end{tabular}
\end{center}

\end{notesbox}

\vspace{0.08in}

\begin{practicebox}{Practice: Evaluate --- Key}
\renewcommand{\arraystretch}{1.8}
\begin{tabularx}{\linewidth}{@{} X X X @{}}
  1. $\arcsin(1) = $ \ans{$\dfrac{\pi}{2}$}
  &
  2. $\arccos\!\left(\dfrac{\sqrt{3}}{2}\right) = $ \ans{$\dfrac{\pi}{6}$}
  &
  3. $\arctan(-1) = $ \ans{$-\dfrac{\pi}{4}$} \\[4pt]
  4. $\sin^{-1}\!\left(-\dfrac{1}{2}\right) = $ \ans{$-\dfrac{\pi}{6}$}
  &
  5. $\cos^{-1}\!\left(-\dfrac{\sqrt{2}}{2}\right) = $ \ans{$\dfrac{3\pi}{4}$}
  &
  6. $\tan^{-1}(0) = $ \ans{$0$}
\end{tabularx}
\end{practicebox}

\end{document}
